\documentclass[11pt,a4paper]{article}

\usepackage[utf8]{inputenc}
\usepackage[T1]{fontenc}
\usepackage[spanish]{babel}
\usepackage{geometry}
\usepackage{array}
\usepackage{longtable}
\usepackage{enumitem}
\usepackage{hyperref}
\usepackage{xcolor}
\usepackage{fancyhdr}
\usepackage{titlesec}

\geometry{margin=2.5cm}
\hypersetup{
    colorlinks=true,
    linkcolor=blue!60!black,
    urlcolor=blue!60!black,
    citecolor=blue!60!black,
    pdftitle={Documentación técnica y funcional de Papeados},
    pdfauthor={Equipo Papeados}
}

\pagestyle{fancy}
\fancyhf{}
\lhead{Papeados}
\rhead{Documentación del proyecto}
\cfoot{\thepage}

\titleformat{\section}{\Large\bfseries}{\thesection.}{0.5em}{}
\titleformat{\subsection}{\large\bfseries}{\thesubsection.}{0.5em}{}

\title{\textbf{Papeados}\\[0.3em] Documentación funcional y técnica}
\author{Equipo de desarrollo de Papeados}
\date{21 de abril de 2026}

\begin{document}

\maketitle

\section*{Autores del proyecto}
\begin{itemize}[leftmargin=*, itemsep=0.6em]
    \item \textbf{Jonatan Godinez Flores} \\
    Rol: Programador de Red \\
    Responsabilidades: Networking, sincronización, lobby \\
    GitHub: \href{https://github.com/}{GitHub}

    \item \textbf{Axel Espinosa} \\
    Rol: Programador de Lógica de Juego \\
    Responsabilidades: Mecánicas, colisiones, estados del juego \\
    GitHub: \href{https://github.com/InozaAki}{InozaAki}

    \item \textbf{Bryan Julián Mendez Ambriz} \\
    Rol: Diseñador UI/UX y QA \\
    Responsabilidades: Interfaces, arte, testing, coordinación \\
    GitHub: \href{https://github.com/Bjma1507}{Bjma1507}
\end{itemize}

\begin{abstract}
Papeados es un videojuego multijugador en red inspirado en el clásico \emph{hot potato}. El objetivo es evitar quedarse con la papa bomba cuando el temporizador llega a cero. Esta documentación describe la estructura general del proyecto, sus sistemas principales, el flujo de juego y la organización de carpetas a partir del código actual del repositorio.
\end{abstract}

\clearpage
\tableofcontents
\clearpage

\section{Descripción general}
Papeados es un juego 2D para entre 2 y 6 jugadores donde una papa explosiva se transfiere entre jugadores mediante colisión. Quien conserva la papa cuando termina el conteo es eliminado y la partida continúa con una nueva ronda hasta que quede un único jugador o se alcance la condición de victoria.

El proyecto está desarrollado en Godot con GDScript y organiza su lógica en módulos separados para facilitar el mantenimiento:
\begin{itemize}[nosep]
    \item Menú principal y lobby multijugador.
    \item Gestión central de la partida.
    \item Jugadores, movimiento y sincronización en red.
    \item Papa explosiva, rondas, puntuación y eliminación.
    \item Interfaz de usuario, escenas finales y gestión de mundo.
\end{itemize}

\section{Objetivo del juego}
El diseño base del juego responde a estas reglas:
\begin{itemize}[nosep]
    \item La partida requiere al menos dos jugadores.
    \item Al iniciar una ronda, un jugador recibe la papa.
    \item El jugador con la papa debe tocar a otro para transferírsela.
    \item Si el temporizador termina con un jugador portando la papa, ese jugador queda eliminado.
    \item El ciclo se repite hasta que un jugador alcance la condición de victoria definida por rondas ganadas.
\end{itemize}

\section{Arquitectura general}
El proyecto separa la lógica en tres niveles:

\subsection{Capa de interfaz}
Incluye el menú principal, el lobby, las pantallas de ganador y perdedor, además de la UI de juego.

\subsection{Capa de juego}
Contiene la lógica de movimiento del jugador, el sistema de papa explosiva, el control de rondas, el sistema de puntuación y la gestión del estado de juego.

\subsection{Capa de red}
Sincroniza jugadores, estados, puntuación y eventos importantes entre servidor y clientes usando RPC y autoridad de red.

\section{Flujo de ejecución}
El recorrido principal de una partida es el siguiente:
\begin{enumerate}[nosep]
    \item El jugador entra al menú principal.
    \item Se crea una sala o se une a una existente desde el lobby.
    \item El host inicia la partida cuando hay suficientes jugadores.
    \item Se carga la escena principal del juego.
    \item \texttt{GameManager} inicializa el estado de red y arranca la primera ronda.
    \item \texttt{RoundManager} prepara el escenario, respawnea jugadores y emite el inicio de ronda.
    \item \texttt{PotatoManager} coloca la papa en un jugador aleatorio vivo.
    \item Los jugadores se mueven y transfieren la papa por colisión.
    \item Cuando ocurre la explosión, se eliminan los afectados y se actualiza el marcador.
    \item Si un jugador alcanza la meta de rondas, se muestran las pantallas finales.
\end{enumerate}

\section{Módulos principales}

\subsection{\texttt{GameManager}}
El gestor central de la partida coordina la interacción entre todos los subsistemas.
\begin{itemize}[nosep]
    \item Decide si la instancia actúa como servidor o cliente.
    \item Conecta señales de papa, rondas y estado de juego.
    \item Spawnea jugadores al entrar nuevos peers.
    \item Maneja desconexiones, reinicio de partida y anuncios de ganador.
    \item Coordina la música del juego y el arranque de la primera ronda.
\end{itemize}

\subsection{\texttt{PlayerManager}}
Administra la vida de los jugadores en la escena.
\begin{itemize}[nosep]
    \item Instancia jugadores con autoridad de red por peer.
    \item Guarda posiciones de aparición rotativas.
    \item Respawnea, elimina y marca jugadores como vivos o muertos.
    \item Provee consultas como obtener jugadores vivos, posiciones o el jugador aleatorio disponible.
\end{itemize}

\subsection{\texttt{PotatoManager}}
Controla la papa explosiva y su ciclo de vida.
\begin{itemize}[nosep]
    \item Spawnea la papa en un jugador vivo.
    \item Transfiere la papa entre jugadores cuando el servidor lo ordena.
    \item Maneja el temporizador de aparición automática.
    \item Recibe el evento de explosión y emite la lista de jugadores afectados.
    \item Gestiona texto flotante de retroalimentación visual.
\end{itemize}

\subsection{\texttt{RoundManager}}
Organiza la secuencia de rondas.
\begin{itemize}[nosep]
    \item Inicia rondas y reestablece estados de supervivencia.
    \item Determina cuándo una ronda termina.
    \item Suma puntos al superviviente cuando corresponde.
    \item Verifica si alguien alcanzó la cantidad de rondas necesarias para ganar.
    \item Envía a las pantallas de ganador o perdedor al cerrar la partida.
\end{itemize}

\subsection{\texttt{ScoreManager}}
Lleva el registro de puntos por jugador.
\begin{itemize}[nosep]
    \item Inicializa el marcador con los peers activos.
    \item Registra nuevos jugadores.
    \item Aumenta la puntuación del superviviente.
    \item Sincroniza los cambios en clientes mediante RPC.
    \item Permite consultar puntos individuales o el mapa completo de puntuaciones.
\end{itemize}

\subsection{\texttt{GameStateMachine}}
Define el estado lógico global de la partida.
\begin{itemize}[nosep]
    \item Estados: \texttt{WAITING}, \texttt{STARTING}, \texttt{IN\_ROUND}, \texttt{ROUND\_OVER} y \texttt{GAME\_OVER}.
    \item Notifica transiciones a través de señales.
    \item Permite reaccionar de forma centralizada a inicio de juego, inicio de ronda, fin de ronda y game over.
\end{itemize}

\subsection{\texttt{Player}}
El script del jugador concentra el comportamiento de movimiento y red.
\begin{itemize}[nosep]
    \item Maneja gravedad, salto, dash y movimiento horizontal.
    \item Reproduce animaciones según el estado del personaje.
    \item Sincroniza posición, velocidad y orientación con la red.
    \item Permite knockback por colisión.
    \item Gestiona la transferencia de papa cuando un jugador entra en contacto con otro.
\end{itemize}

\subsection{\texttt{PlayerSprite}}
Genera el efecto visual de estela o \emph{ghost} durante el dash.
\begin{itemize}[nosep]
    \item Crea una caché de sprites duplicados para reutilizarlos.
    \item Activa el efecto solo cuando el jugador está en dash.
    \item Anima el desvanecimiento del sprite para simular velocidad.
\end{itemize}

\section{Sistema multijugador}
La base online utiliza un modelo cliente-servidor donde el servidor es autoridad para los eventos importantes.

\subsection{Autoridad}
\begin{itemize}[nosep]
    \item Cada jugador controla su propio personaje local.
    \item El servidor decide la creación de jugadores, transferencia de papa, eliminación y puntuación.
    \item Los clientes reciben los cambios por RPC y actualizan su vista local.
\end{itemize}

\subsection{Sincronización}
Los datos que se replican con más frecuencia son:
\begin{itemize}[nosep]
    \item Posición y velocidad del jugador.
    \item Cambios de estado de la partida.
    \item Estado de la papa y su transferencia.
    \item Marcador y resultados de ronda.
\end{itemize}

\subsection{Conexión y lobby}
El lobby permite crear o unirse a una sala, mostrar participantes y validar que haya suficientes jugadores antes de iniciar la partida. Una vez que el host empieza, la escena principal se carga para todos los participantes.

\section{Interfaz y escenas}
El proyecto separa la presentación en escenas distintas:
\begin{itemize}[nosep]
    \item Menú principal para crear o unirse a una partida.
    \item Lobby multijugador con lista de jugadores y control de inicio.
    \item Escena principal de juego.
    \item Pantalla de ganador.
    \item Pantalla de perdedor.
\end{itemize}

También existe una capa de UI para HUD, textos flotantes y elementos auxiliares. La escena de fondo y los elementos visuales refuerzan la identidad del proyecto con una estética arcade colorida.

\section{Mundo y escenarios}
La carpeta \texttt{world} agrupa elementos del entorno:
\begin{itemize}[nosep]
    \item Generación de plataformas.
    \item Bordes del mundo y límites del mapa.
    \item Cambios de gravedad.
    \item Plataformas especiales como lanzadores o interruptores.
\end{itemize}

Este bloque permite extender el diseño de mapas sin tocar la lógica de jugador o de papa.

\section{Estructura de carpetas}
La siguiente organización resume el contenido actual del proyecto:

\begin{longtable}{|>{\raggedright\arraybackslash}p{3.3cm}|p{10cm}|}
\hline
\textbf{Carpeta} & \textbf{Propósito} \\
\hline
\endhead
\texttt{assets/} & Recursos artísticos y de audio del juego. \\
\hline
\texttt{core/} & Sistemas centrales: estado de aplicación, estado de juego, red, validación y manager principal. \\
\hline
\texttt{gameplay/} & Reglas de juego, rondas, puntaje, jugadores y papa. \\
\hline
\texttt{multiplayer/} & Lobby y flujo de conexión. \\
\hline
\texttt{player/} & Lógica y escena del jugador. \\
\hline
\texttt{potato/} & Lógica y escena de la papa explosiva. \\
\hline
\texttt{scenes/} & Menús, pantallas finales, administración de escenas y música. \\
\hline
\texttt{ui/} & Interfaz, HUD, menús y textos flotantes. \\
\hline
\texttt{world/} & Generación del escenario y elementos del mapa. \\
\hline
\end{longtable}

\section{Tecnologías y herramientas}
\begin{itemize}[nosep]
    \item Motor: Godot.
    \item Lenguaje: GDScript.
    \item Control de versiones: Git y GitHub.
    \item Diseño visual: Figma y Piskel.
    \item Comunicación y coordinación: Discord y Notion.
\end{itemize}

\section{Estado actual del proyecto}
Con base en el estado del repositorio, el proyecto se encuentra en desarrollo activo. La base funcional ya incluye:
\begin{itemize}[nosep]
    \item Menú principal y lobby multijugador.
    \item Gestión central de partida.
    \item Spawns de jugador y eliminación.
    \item Transferencia de la papa y explosión.
    \item Sistema de rondas y marcador.
    \item Pantallas de victoria y derrota.
\end{itemize}

Quedan por consolidar o pulir aspectos como la experiencia visual final, la optimización de red, el balance del gameplay y el empaquetado final de documentación.

\section{Plan para la implementación de la conexión en red en la siguiente fase}
Aunque el proyecto ya cuenta con una implementación funcional de conexión en red, para efectos del entregable se plantea un plan de implementación y consolidación que modela cómo habría sido la fase siguiente de desarrollo y mejora.

\subsection{Objetivo general}
Diseñar, estabilizar y validar la arquitectura de red del juego para garantizar partidas multijugador fluidas, sincronización consistente entre clientes y servidor, y una base escalable para nuevas mecánicas.

\subsection{Objetivos específicos}
\begin{itemize}[nosep]
    \item Definir y documentar el modelo de autoridad (servidor autoritativo y clientes replicados).
    \item Estandarizar los RPC críticos de juego (spawn, estado de ronda, transferencia de papa, puntuación y fin de partida).
    \item Reducir desincronizaciones de movimiento mediante interpolación y umbrales de actualización.
    \item Fortalecer la tolerancia a desconexiones y reconexiones en mitad de partida.
    \item Establecer pruebas de estrés para medir estabilidad y latencia percibida.
\end{itemize}

\subsection{Plan por etapas}
\begin{enumerate}[nosep]
    \item \textbf{Etapa 1: Diseño técnico de red} \\
    Definición de eventos sincronizados, estructura de mensajes RPC, reglas de autoridad y diagrama de flujo de conexión lobby $\rightarrow$ partida.

    \item \textbf{Etapa 2: Implementación base} \\
    Integración de creación/unión de sala, spawn de jugadores en servidor, replicación de estado inicial y sincronización de transición de escenas.

    \item \textbf{Etapa 3: Sincronización de gameplay} \\
    Replicación de movimiento, orientación, estados de ronda, posesión de papa, explosión y eliminación, con validaciones del lado servidor.

    \item \textbf{Etapa 4: Robustez y recuperación} \\
    Manejo de peer desconectado, reasignación de flujo de ronda, recuperación de estado para clientes tardíos y reinicio seguro de partida.

    \item \textbf{Etapa 5: Optimización y QA} \\
    Ajuste de tasa de actualización, reducción de tráfico innecesario, pruebas con múltiples clientes y corrección de jitter o regresiones.
\end{enumerate}

\subsection{Cronograma tentativo}
\begin{itemize}[nosep]
    \item Semana 1: Diseño técnico y contrato de mensajes de red.
    \item Semana 2: Implementación base de conexión, sala y sincronización inicial.
    \item Semana 3: Integración completa del gameplay sincronizado.
    \item Semana 4: Pruebas de carga, manejo de fallos y optimización final.
\end{itemize}

\subsection{Riesgos técnicos y mitigación}
\begin{itemize}[nosep]
    \item \textbf{Desincronización de estados}: Mitigar con servidor autoritativo y eventos confiables para cambios críticos.
    \item \textbf{Latencia percibida en movimiento}: Mitigar con interpolación en clientes y umbrales de envío.
    \item \textbf{Errores por desconexión}: Mitigar con limpieza controlada de peers y reanudación de ronda.
    \item \textbf{Regresiones al escalar}: Mitigar con pruebas por escenario y checklist de validación por release.
\end{itemize}

\subsection{Entregables de la fase}
\begin{itemize}[nosep]
    \item Documento técnico de arquitectura de red y decisiones de diseño.
    \item Matriz de RPC con propósito, confiabilidad y origen de autoridad.
    \item Registro de pruebas multijugador (funcionales y de estrés).
    \item Informe de optimización de red con métricas de estabilidad.
\end{itemize}

\section{Problemas encontrados y reestructuración}
Durante las primeras etapas del desarrollo, el proyecto se construyó con una estructura más monolítica. En ese momento parecía la opción más rápida y sencilla para avanzar con un prototipo funcional, ya que concentrar gran parte de la lógica en pocos scripts facilitaba implementar mecánicas básicas en poco tiempo.

Sin embargo, conforme se agregaron nuevas funcionalidades (multijugador, control de rondas, sincronización de estado, UI y reglas de eliminación), esa aproximación comenzó a volverse insostenible. Los principales problemas fueron:
\begin{itemize}[nosep]
    \item Alto acoplamiento entre sistemas, donde un cambio en una mecánica afectaba partes no relacionadas.
    \item Dificultad para depurar errores, especialmente en lógica de red y flujo de partida.
    \item Menor claridad del código para trabajo colaborativo, al no existir límites claros de responsabilidad.
    \item Mayor riesgo de regresiones al extender funcionalidades o corregir bugs.
\end{itemize}

Como resultado, se realizó una reestructuración hacia la arquitectura modular actual, separando responsabilidades por managers y sistemas especializados (jugadores, papa, rondas, puntuación, estados y red). Esta transición mejoró significativamente la mantenibilidad del proyecto, permitió escalar el desarrollo de forma más segura y simplificó la incorporación de nuevas mecánicas.

\section{Conclusión}
Papeados es un proyecto académico con una base modular clara. La separación entre gestores de juego, red, jugador, papa y UI facilita ampliar reglas, agregar mecánicas nuevas y mantener el código. Esta documentación puede servir como referencia técnica para el equipo y como base para entregas futuras.

\end{document}